\section{怎样进行审题}


中学生写作文跟作家进行创作的不同特点之一，就是要按题作文，因为作文只是作为学习语文的一种训练手段，它是要按照教学计划有步骤地进行的，以便循序渐进地培养写作能力。既然不论平时训练还是考试，作文都要有题目，因此在作文时，就得琢磨透题目的含义和要求，明确题目要求写什么内容和应该用什么形式来表达，然后才好下笔。这个琢磨题意的过程叫做审题。

审题的能力就是正确地理解和分析题目要求的能力。审题训练是作文训练的第一步，也是关键的一步。只有题意审准了，审对了，写出的作文才能切题，符合要求；如果题意领会偏了，领会错了，写出的作文就会离题，甚至文不对题。

那么，我们怎样才能审清题意，为写好作文打下坚实的基础呢？

就审题的步骤来说，大致有二：一是弄清题目对作文的内容和形式究竟提出了哪些要求；二是发掘题目除了限定的范围或条件之外，还有哪些不曾限制而可供充分发挥的领域，以便打开思路，写好文章。

\textbf{1. 弄清题目对作文的内容和形式究竟提出了哪些要求。}

从内容上看，作文题或是规定了写人、记事的范围，或是规定了时间、空间的范围，或是规定了人称、数量的范围；有的规定了论述或批驳的问题，有的则规定了论证的中心，还有的规定了说明的对象，等等。凡是题目规定的范围和要求，作文时必须遵守，否则就不合要求。

例如，《我在成长过程中的二三事》这个作文题，它要求以“我”作文章的中心人物，规定了人称；它还要求记“事”，而且规定了数量“二三”（即是说所记之事的数量不得少于“二”，也不得多于“三”）；另外还规定了“二三事”的时间和空间的范围“在成长过程中”。很明显，这篇作文应该写成写人记事的比较复杂的记叙文。如果以“我们”或别人作中心人物，或者只写出一件事或三件事以上，或者所写的内容与“在成长过程中”无关，或者主要不是记“事”而是写在成长过程中有什么心得体会，就算是离题了。

再如，《童年忆趣》这个作文题，规定我们要“忆”，即要求写回忆中的事；它还规定了回忆的时间局限在“童年”时期；它又规定所回忆的事应具有“趣”的性质。很明显，这篇作文也应该写成写人记事的记叙文。如果我们不写回忆中的事，而写现实中的事，或者不写童年时代的回忆，而写少年时代、青年时代的事；或者写的事并不具有“趣”的性质，而是具有悲痛、严肃或惊心动魄的性质，就算是离题了。

又如，《求知无坦途》这个作文题，是用一个否定判断的句式出现，说的是一个观点，要我们充分论证它是一种正确的主张或见解，很明显，它要求我们写一篇以立论为主的议论文。它规定论证的主要对象是“求知”，即追求真知，并且规定论证这个“求知”的过程“无坦途”，即没有平坦轻松的大道可走，必须沿着崎岖艰难的小路前进。如果我们不是论证“求知”，而是论证品德修养、锻炼身体或者社会主义建设、社会主义革命；如果不是论证求知“无坦途”，而是论证什么精益求精啦，循序渐进啦，摸清规律啦，锲而不舍啦，就都算离题了。

又如，《张华牺牲得不值得吗？》这个作文题，是用一个反问句子的形式出现，它的实际含义当然是说张华这位舍己救人的新型大学生“为人民而死是比泰山还要重的”，但是，句子既以反问形式出现，则明确要求我们批驳那种认为“张华牺牲不值得”的错误观点，即要求我们写一篇以驳论为主的议论文。如果我们不是把张华作为主要的评价对象，而把刘胡兰、黄继光或雷锋作为主要的评价对象，或者批驳的内容不是“牺牲得不值得”，而是去批驳“吃喝玩乐腐朽人生观”，去揭露资本主义社会尔虞我诈的反动本质，或者把它写成正面立论的文章，等等，显然都是不合要求的。

从形式上看，作文题一般都已规定了文章的体裁和写法。记叙文的题目，往往带有“传”、“记”、“忆”、\\ 
“记······人”、“记······事”、“······回忆”等一类字眼；议论文的题目，往往带有“议”、“论”、“说”、“评”、“析”、\\ 
“谈”、“批”、“驳”、“斥”、“读”、“有感”、“启示”、“学习······的体会”、“从······想到的”等一类字眼；说明文的题目，则往往带有“说明”、“介绍”、“简说”、“简介”等一类字眼；应用文的题目更加明显，一般都写明“信”、“日记”、“笔记”、“发言稿”、“演讲稿”、“总结”、“纪要”、\\
“调查报告”、“通知”、“计划”等字眼。如果作文题上面没有明显的识别标志，那就要从题目所涉及的事物来辨别。题目涉及的要是具体的人、物、事、地、时，则一般要求写记叙文，如《净友》、《苦战》、《凡人小事》、《难忘的一课》、《彩色的夜》、《考试结束以后》、《马路见闻》、《我所经历的一段有意义的生活》、《故乡重逢》，等等；题目涉及的要是较抽象的道理、主张、问题、口号等，则一般要求写议论文，如《为振兴中华而刻苦学习》、《个人要服从集体》、《千里之行，始于足下》、《不知足才常乐》、《心底无私天地宽》、《万丈高楼平地起》、《想和干》、《有了文凭就有一切吗？》，等等。题目涉及的要是种客观事物、一件工艺品、一门技术、一种方法等，则一般要求写说明文，如《南京长江大桥》、《蜘蛛》、《冰灯》、《心灵的窗户 --- 眼睛》、《菊花》、《怎样查字典》、《锻炼记忆力的方法》，等等。

总之，拿到一个作文题，必须全面地（而不是片面地）、过细地（而不是粗率地）理解题目中每个词的含义，弄清各个词之间的内在联系，从而弄清题目对文章的内容、文体和写法等作了哪些要求和规定；如果题目中供给一些材料，还得仔细研究这些材料，确定其中心意思和关键词句，再根据材料的中心意思，紧扣所给的作文题，下笔作文。近几年来的高考作文题，多数以供给限定的材料，再进行命题的形式出现，其审题法与一般的命题作文审题法在基本原则上是—致的。

\textbf{2. 发掘题目中不曾限制而可供充分发挥的领域。}

弄清题目所给的限制，按照这些规定作文，而不能随心所欲地改变题目规定的要求，这是审题的第一步。但是不能就此止步，而被这些限制束缚了思路。应该再进一步思考：在题目限定的范围内，还有哪些没加限制而可供发挥的地方。这比起第一步来，显得更加重要，因为第一步工作一般来说比较容易做到，但有不少人往往在下笔作文时被这些限制捆住了手脚。只有迈出第二步，才能广开思路，结合自己的特长和见闻，写出丰富多彩、新颖动人的文章来。

例如，记叙文题目总是在时、地、人、事等方面作出一定的规定或给一些具体材料，但是又不能在一个题目中把各种限制条件都包括在内；即使在有所限制的范围内，也不可能把具体写哪一方面限制得太死。那么，我们就可以从那限制不了的部分来开拓思路。如果写《班主任》，题目既已规定了要写的人是“班主任”，当然不能把其他的任教老师或者同学作为中心人物来写。但是，它并没规定是谁的班主任，也没规定是哪个学校，哪个班级的班主任，更没有规定必须写他的外貌、行为、内心活动，或者怎样做班主任的，他与同学的关系又怎么样，等等，这就使我们大有发挥、联想的余地，刘心武同志甚至以这个题目写了一篇著名的小说。如果是写《我的班主任》或《我最难忘的班主任》等等，限制多起来了，就只能写给自己任教的班主任，或者是写自己班主任中留给自己印象最深的一位。但是，题目并没有限制是写小学时代还是中学时代的哪一位班主任，也没有限制写他的一件事还是几件事，是写他教学上的事、工作上的事、关怀学生的事，还是生活中感人的事，更没有限制写他在课堂上的事、校内的事、家内的事还是在家访中的事、路上的事、医院中的事，等等。因此，我们只要结合自己的亲身经历和感受，从中选取使自己永难忘记的一件或几件事情，来表现某位班主任的“园丁”形象和高尚品德，就能写出既切题意义多种多样、新颖动人的作文来。这就是说，对题目给予的规定，所给的材料，既要确切理解，融会贯通，又不可太死板，被它束缚住，而要从它不曾限定的地方打开思路，做到既扣紧题意，又挥洒自如。

对于议论文的题目，我们也可以大致这么办。议论文的题目，一般要提出某一个比较抽象的概念或观点（即使题目只提了一个问题或只供给一点材料，经过分析，也可以提炼成某一概念或观点）。我们可以把这个概念或观点变成“是什么”、“为什么”和“怎么样”等问题来展开思考。例如一九八二年的高考作文题《先天下之忧而忧，后天下之乐而乐》，它用两个意义相对的短语表达了一个观点，那么这个观点的确切含义是什么？“先”和“后”，“忧”和“乐”的截然相反，而“忧”和“乐”的对象又同是“天下”，这表现了一种什么样的思想境界呢？为什么会这样呢？从古至今能够这样做的是哪些人，他们是什么样的人？与此相反的是哪些人，他们又是什么样的人？第一个写出这句话的宋代大政治家范仲淹与当代的共产主义战士对这句话的理解又有什么不同之处呢？在社会主义时代，我们应该怎样做一个为了“天下”人民的彻底解放而“吃苦在前，享乐在后”的人呢？等等。这些问题既在这个题目论述的范围之内，又是题目没有加以限制的地方，我们可以根据这些问题开拓思路，写出既扣紧题目又议论深刻的文章来。

我们除了懂得审题的这两个基本步骤之外，还应该掌握一定的方法。大致说来，以下几种审题方法值得注意。

\textbf{1. 从文字结构上进行分析。例如：}

①\ （细）\uuline{流}与（大）\uuline{海}（联合结构）

②\ （文明）之\uuline{花}[\ 处处\ ]\uline{开}（主谓结构）

③\ \uline{记}（一件）（有意义）的（小）\uwave{事}（动宾结构）

④\ \uline{写}(在(\ \uuline{百花}\ \uline{争}\ \uwave{春})的日子里)（动补结构）
\newpage

⑤\ （我）的（第一个）（领路）\uuline{人}（偏正结构）

⑥\ $\langle$ 在（大战红五月）的日子里 $\rangle$ （介词结构）

对题目的词序和句法进行分析，可以帮助我们弄清题目字词之间的关系，确定文章的体裁，明确题目限制的范围、内容和数量等，同时也可以帮助我们思考题目没加限制的领域。如“细流与大海”这个作文题中的“细流”与“大海”之间是并列关系，题目要求将二者之间的辩证关系论述清楚，着眼点是在“小”与“大”、“流”与“海”的相反相成上；而“细流”和“大海”之间的联系和比较能表达出一个什么样的观点？与“细流”相类似的事物是哪些？与“大海”相类似的事物是哪些？它们之间的关系又是如何？为什么？我们从它们之间的联系和比较中能够得到什么样的启示？这些问题既在这个题目考虑范围之内，又是题目没有加以限制的地方，我们可以从这些问题中打开思路。

\textbf{2. 抓住关键性的词语。}

一个题目中的每个词不可能一样的重要，一般来说，题目中的谓语和修饰语往往是题意的重点之所在。有人把这道意的重点叫做“题眼”，因为只有找出了重点，下笔时才能把握文章的中心，才好适当地选材布局，突出主题思想。如上面所列的六个题目，重点所在分别是“细”与“大”，“处处开”，“有意义”，“百花争春”，“领路”，“大战红五月”。再如《歌声又响起来了》，重点词是“又”字，应着重写出“又”一次响起来的情景和意义；《课堂新风》，重点词是“新”字，“新”在哪里，怎样“新”法，为什么“新”，就成了文章的中心所在。

\textbf{3.补足省略部分。}
有些题目，只是标明一个线索，暗含了某些省略的部分，审题时就要开动脑筋，多问几个“谁”、“是什么”、“为什么”、“怎么样”等。把题目中省略的部分补足，然后据此作文。如《在红旗下》这个题目，既无主语，又无谓语，只是个介词结构，我们可以问：“谁”在红旗下？在红旗下“怎么样”，“干什么”？如果把这个题目补足理解成《我在红旗下成长》、《他在红旗下前进》等，就适宜写记叙文；如果补足理解成《我在红旗下想到的》，就适宜写成以抒情为主的散文。再如《苦战》，可以补足理解成《我们在苦战》、《苦战在平凡岗位上》等，这适宜写记叙文；也可以补足理解成《论苦战》、《为实现四化而苦战》、《从苦战谈起》等，则适宜写成议论文。这就是说，象《苦战》、《园丁》、《在大战红五月的日子里》等标明事件、人物、时间因素，本应写成记叙文的题目，如果规定要我们写成议论文的话，就要运用补足省略成分的方法，使其成为所规定的文体的题目。

\textbf{4. 揣摩题目寓意。}

有些作文题目，既有字面的具体意义，又有某种比喻的、象征的抽象含义。例如《细流与大海》，《我的第一个领路人》；又如《路标》、《春雨》、《晨》、《路》、《铺路石》、《火炬》、《要知松高洁，待到雪化时》、《不到长城非好汉》等等。对这一类题目，我们不仅要把它们的字面意思与自己经历过的某些事联系起来，而且要体会它们人为地赋予的、不是这个词语所固有的意义。由于你赋予它的象征意义不同，题目的范围和要求也就不同，文体也会随之不同。那就要根据自己情况，扬长避短，选择自己擅长的文体和角度落笔写文了。

\textbf{5. 拟出若干个相似的题目加以比较。}

这种方法可以使我们明白：相类似的题目之间总是有细微的差别，而这点差别正是一个题目独具的范围和要求。例如：《记一位优秀团员》、《记一位三好生》、《三好学生×××的事迹鼓舞着我》。前两题的相同点都是要记一个好青年。但前一题要从“团员”的角度写出他“优秀”的事迹，他可以是学生，也可以是工人或是战士；后一题则要从“学生”的角度写出他“三好”的事迹；至于第三道题的重点并不在于写三好学生的事迹，而要侧重写出“鼓舞着我”，着眼于“我”在“鼓舞”下的进步。再如：《水滴自能石穿》、《水滴自能石穿吗？》《水滴石穿辨析》。前两题的共同点，都是要从水滴石穿这个自然现象的象征性意义论说一个道理。但是前一题是用肯定的语气表述，要求正面地论证从这个自然现象得出的一个真理：只要目标如一，坚持到底，就能成功；后一题则是用反问的语气表述，要求从反面批驳那种只想等待必然的结果而不尽主观努力的懒汉思想，论说“一万年太久，只争朝夕”的积极进取精神。至于第三个题目着重在“辨析”，要求从正面、反面等不同角度去辨别、分析这个自然现象给予人们的各种启示，论说既要目标如一、坚持到底，又要发挥主观努力、积极进取的思想。经过这样的比较，各题的细微差别就被找了出来，题意也就审清了。
\newpage